\section{Production-hardening lessons}\label{sec:lessons}

The thesis of \S\ref{sec:intro} has a second, less statistical face: several defects in this system
were invisible to its own test suite precisely because the path the tests exercised was not the path
production takes. Each lesson below was found by an audit or an experiment rather than by a failing
test, was measured before it was fixed, and is recorded in the design log
(Appendix~\ref{app:deviations}).

\paragraph{L1: the shared resource nobody declared concurrency-unsafe.}
The replicated design asserts that execution is deterministic and auditable. Versioned partition
creation read the current maximum version and then dropped, created and registered a table in
separate statements with no lock; two concurrent callers for the same partition computed the same
next version and raced. It surfaced as an intermittent integration failure in continuous
integration, and was root-caused rather than retried: a recovery-agent replay triggered the same
DAG as an unrelated test. The fix serialises the read-modify-write with a database advisory lock.
\emph{Lesson:} an agent that can trigger real work turns every shared resource into a concurrency
surface that a single-actor test never touches.

\paragraph{L2: an action could be lost from the audit trail.}
The audit row was written \emph{after} execution. A crash, out-of-memory kill or database blip in
that window meant the action had really happened with no record, and because the control loop
swallows per-tick exceptions the failure was silent. For a system whose central claim is that every
action is policy-gated and auditable, an unauditable executed action falsifies the claim. The
intent row, with the policy verdict already known, is now committed \emph{before} execution and
updated with the outcome.

\paragraph{L3: the detector that had no callers.}
The table that recovery-time statistics and decision scoring depend on was written only by the
fault injector. The monitoring agent updated an existing row's detection time but never created one,
and a complete, unit-tested deterministic anomaly detector had no caller at all. In a deployment
without the injector, a genuinely detected incident reached the audit table but never became an
incident: production could not have reproduced the experiment's metrics. Found while planning
features, by reading, not by a test.

\paragraph{L4: two taxonomies that never met.}
Decision-correctness is scored against an accepted-mitigation set keyed by the \emph{injected
scenario} names. Applying it to live incidents, whose fault types have different names, would have
scored every real incident ``incorrect by construction''. A mapping for live fault types was added
before the metric was exposed.

\paragraph{L5: hermeticity by enumeration does not scale.}
A recurring unit-test failure class --- a test forgetting to mock one module's database handle,
turning a millisecond test into a thirty-second pool timeout --- was fixed instance by instance four
times before the mechanism was addressed: every such call funnels through a single connection-pool
accessor, which is now patched to fail fast for all tests. Mutation-testing the guard (removing one
load-bearing mock) failed six tests in seconds with a message naming the cause. \emph{Lesson:} when
the fourth instance appears, fix the mechanism.

\paragraph{L6: what the adversarial corpus found.}
Grading a generated corpus against an independent oracle found three things (numbers in
\S\ref{sec:results}). (i) The cost policy prices a scaling target against the budget but says
nothing about its floor: zero workers, which halts ingestion, or a negative ``scale-down'', was
budget-legal, and non-numeric targets raised inside the gate before the audit row was written, so a
rejected proposal left no record. The fix is in the contract layer: scaling targets must be integers
at least one, and violations take the ordinary invalid-output path. (ii) Against a newer policy-engine
release than the one the project pins, some float budget comparisons evaluated as within budget
when they were not; the pinned release was correct. We report the observation only: we did not
root-cause it or report it upstream. The consequence stands regardless: the pin is load-bearing, and
the corpus is now an integration test that would catch a version bump that changes the semantics.
(iii) Once the contract was fixed, the policy agreed with the independent specification on every
graded case: the containment claim is now a measured property with an interval, not an assertion.

\paragraph{L7: the pilot found a dead model.}
Before the live campaign, an eight-run pilot at the real timings was run to measure cost and time.
Within minutes it found that the provider had retired the small model configured for the monitoring
agent (HTTP 410, end of life). The monitoring agent is the detector for every non-baseline
configuration, so a full campaign would have produced a matrix in which the detector silently
degraded to a no-op and the numbers would have measured provider retirement. The supervisor's
degraded-streak guard would have caught it after several wasted runs; the pilot caught it at the
first call. \emph{Lesson:} the cheapest validity check in a long, paid experiment is a short one at
the real settings, and a listed model is not a usable model.
